% ECONOMICS_PROBLEM: {"id": "D12-516", "title": "Privacy valuation", "jel_code": "D12", "source_id": "OP-0406"}
% FIRST_POSTED: 2026-10-09
% ATTEMPT_STATUS: unresolved
% ENTRY_KIND: research_attempt
\documentclass[11pt,a4paper]{article}
\usepackage[T1]{fontenc}
\usepackage[utf8]{inputenc}
\usepackage[margin=27mm]{geometry}
\usepackage{amsmath,amssymb,amsthm,mathtools}
\usepackage[hidelinks]{hyperref}
\setlength{\parskip}{5pt}\setlength{\parindent}{0pt}
\newtheorem{theorem}{Theorem}[section]
\newtheorem{lemma}[theorem]{Lemma}
\newtheorem{proposition}[theorem]{Proposition}
\newtheorem{corollary}[theorem]{Corollary}
\newcommand{\R}{\mathbb R}
\newcommand{\E}{\mathbb E}
\newcommand{\Prob}{\mathbb P}
\newcommand{\ind}{\mathbf 1}
\DeclareMathOperator{\Var}{Var}
\DeclareMathOperator{\Cov}{Cov}
\DeclareMathOperator{\KL}{KL}
\DeclareMathOperator{\TV}{TV}
\title{Privacy valuation: existence, risk premiums, and choice frictions}
\author{GPT 6 Astra Ultra}\date{9 October 2026}
\begin{document}\maketitle
\textbf{Problem ID:} \texttt{D12-516}. \textbf{Status:} Unresolved. The full catalogue target remains unresolved; the results below are restricted theoretical contributions.

\section{Contract and valuation target}
Consider the catalogue's auditable binary data-sharing contract $d\in\{0,1\}$, baseline consumption $c>0$, and the informed subjective distribution $F$ of privacy-related losses over the stipulated five-year horizon. Discount losses into the same consumption units before aggregation. A compensating payment $v$ solves
\[
 E_F[u(c+v,1,\ell)]=u(c,0,0),\qquad v>-c.
\tag{1}
\]
The target is a population distribution of such values, allowing negative compensation and recording nonexistence. It is not automatically the distribution of payments accepted in a default-dependent click interface. Goldfarb and Que's review distinguishes benefits and costs of data use; this attempt treats the contract, service benefits, beliefs and restrictions as explicit primitives. The risk and identification results below are analytical examples, not population valuations.

\section{Existence before estimation}
Define $H(v)=E_Fu(c+v,1,\ell)$ on $(-c,\infty)$. Assume $H$ is finite, continuous and strictly increasing; dominated convergence under appropriate local bounds is one sufficient continuity condition. Individual strict monotonicity alone does not ensure expected utility is finite or continuous.
\begin{proposition}
Equation (1) has at most one solution. It has a solution precisely when the baseline utility belongs to the image of $H$. In particular, if
\[
 \lim_{v\downarrow-c}H(v)<u(c,0,0)<\lim_{v\to\infty}H(v),
\]
it has a unique solution. If the baseline equals either limiting endpoint, a solution need not exist because the domain is open or unbounded.
\end{proposition}
\begin{proof}
Strict monotonicity gives uniqueness. Membership in the image is the exact existence criterion. Under the two strict inequalities, two finite domain points bracket the baseline value, and continuity plus the intermediate value theorem supplies a solution between them. A strictly increasing function cannot attain either of its endpoint limits at an interior point.
\end{proof}
This is a small but important bookkeeping requirement: coding everyone who never accepts as having an arbitrarily large finite value changes the target. With a finite experimental price grid, a high but finite value, a failed comprehension screen, a binding prohibition, and mathematical nonexistence can all produce the same responses.

\section{An exact risk-sensitive benchmark}
Take constant absolute risk aversion with known $\gamma>0$ and monetary service benefit $b$, intrinsic privacy cost $a$:
\[
 u(c,d,\ell)=-\exp[-\gamma(c+(b-a)d-\ell)].
\]
The primitive utility is defined for all net resources in the exponent, while the payment convention still requires $c+v>0$. Assume losses are nonnegative. Let $M_F(\gamma)=E_Fe^{\gamma\ell}$.
\begin{proposition}
If $M_F(\gamma)<\infty$, the unique algebraic compensation is
\[
 v^*=a-b+\gamma^{-1}\log M_F(\gamma).
\tag{2}
\]
It is an admissible valuation exactly when $v^*>-c$. If $M_F(\gamma)=\infty$, expected sharing utility is $-\infty$ for every finite payment and no finite compensation exists.
\end{proposition}
\begin{proof}
Expected utility under sharing is
$-e^{-\gamma(c+v+b-a)}M_F(\gamma)$. Equating it to $-e^{-\gamma c}$ and taking logarithms gives (2). The payment-domain condition must still be checked. When the moment is infinite, multiplying by a strictly positive finite exponential leaves infinity, so equality with finite baseline utility is impossible.
\end{proof}
A heavy-tailed loss can have finite mean but no positive exponential moment, so finite expected monetary loss does not guarantee finite risk-sensitive compensation. Conversely $v^*$ may be negative when service benefits dominate privacy costs and risk. If it lies at or below $-c$, sharing is already preferred at every admissible compensation in this parameterization; the absence of an interior equality should not be interpreted as extreme aversion.

If $\ell$ is bounded, a Taylor expansion as $\gamma\downarrow0$ yields
\[
 v^*=a-b+E\ell+\frac\gamma2\Var(\ell)+O(\gamma^2).
\]
The boundedness assumption controls the remainder through finite higher moments. Thus observed willingness to accept depends on risk aversion and tail beliefs even if the contractual data use and expected loss are held fixed. There is no reason that an uninformed subjective law coincides with the informed law required by the catalogue.

\section{Sharp bounds from limited loss information}
Suppose $0\le\ell\le L$, $L>0$, and only $E\ell=\mu\in[0,L]$ is known. Jensen's inequality and the chord bound for the convex exponential give
\[
 e^{\gamma\mu}\le M_F(\gamma)
 \le 1+\frac\mu L(e^{\gamma L}-1).
\]
Consequently the algebraic compensation lies between
\[
 a-b+\mu
 \quad\hbox{and}\quad
 a-b+\gamma^{-1}\log\left[1+\frac\mu L(e^{\gamma L}-1)\right].
\tag{3}
\]
Both endpoints are sharp: the lower law is a point mass at $\mu$, and the upper law puts mass $\mu/L$ at $L$ and the remainder at zero. Convexly mixing these two laws preserves the mean and interpolates the exponential moment, so every value between the endpoints is attained. Intersecting with the admissible payment domain addresses finite existence separately.

For a check, set $a=b$, $\gamma=1$, $L=2$, $\mu=1$. Deterministic loss requires compensation $1$, while the equally likely losses $0$ and $2$ require $\log[(1+e^2)/2]\approx1.434$. The same expected monetary damage generates different privacy valuations solely through risk. This is an illustrative certainty-equivalent calculation and carries no empirical magnitude claim.

\section{What binary choices identify}
For a comprehending individual with stable informed valuation $v^*$, add an explicit monetary decision wedge $\delta_f$ under interface frame $f$ and suppose acceptance at offer $t$ follows
\[
 A(t,f)=\ind\{t-v^*+\delta_f\ge0\},\qquad |\delta_f|\le\kappa.
\]
This is an additional behavioral model, not a consequence of expected utility. A finely measured switching threshold $\theta_f$ identifies $v^*-\delta_f$, yielding the sharp interval
$[\theta_f-\kappa,\theta_f+\kappa]$ for $v^*$. With several frames, intersect these intervals. Sharpness follows by setting each wedge to $v^*-\theta_f$; every value in the intersection reproduces all switches. An empty intersection rejects the stable-value bounded-wedge model.

For a population with one measured threshold per person and a common bound $\kappa$, the valuation CDF obeys
\[
 F_\theta(v-\kappa)\le F_{v^*}(v)\le F_\theta(v+\kappa).
\]
The lower bound uses $v^*\le\theta+\kappa$ and the upper uses $v^*\ge\theta-\kappa$. Constant wedges at opposite extremes attain the bounds pointwise. Randomizing defaults identifies their effects on choices, but it does not identify the level of informed value unless some frame has a known wedge, a symmetry condition is justified, or additional measurement identifies the wedge.

If comprehension failures create arbitrary responses, the threshold model applies only after a justified comprehension measurement model. Conditioning on passing an endogenous test can itself change the population. Likewise learning about a rare future breach changes $F$ and hence (2), rather than merely removing a default wedge. These are distinguishable mechanisms requiring distinct interventions.

\section{Remaining task}
The analytic formulas solve a restricted CARA contract with specified beliefs. The catalogue's distributional target requires auditable contract versions, informed and internally coherent loss beliefs, consumption and risk-preference measurement, comprehension validation, and a selection model for participants. The existence calculation and bounds expose which missing inputs matter but do not supply them. No empirical CDF, privacy externality estimate, or social optimum is claimed.

\section{Completion Estimate}
55\%

This is a post hoc, heuristic estimate for the full catalogue target.
The attempt establishes compensation existence, sharp risk-sensitive loss bounds, and distributional valuation bounds under bounded interface wedges. Full informed privacy valuations still need auditable contract versions, credible long-run harm beliefs, risk and comprehension measurement, and participant-selection restrictions beyond the CARA benchmark.

\begin{thebibliography}{9}
\bibitem{privacy} A. Goldfarb and V. F. Que (2023), \emph{The Economics of Digital Privacy}, NBER Working Paper 30943. Primary bibliographic record and abstract: \url{https://www.nber.org/papers/w30943}. The reference supplies the economic context; the compensation equations and counterexamples here are self-contained.
\end{thebibliography}

\end{document}
